\chapter*{Abstract}
\addcontentsline{toc}{chapter}{Abstract (English)}
\markboth{Abstract}{Abstract}

\label{sec:abstract-english}


\textbf{Deep Learning for Two-Hop Relay Communication: A Comparative Study of Classical and Neural Network-Based Strategies}

This thesis presents a comparative study of classical and neural network-based relay strategies for two-hop cooperative communication. Nine relay methods are evaluated: two classical approaches : amplify-and-forward (AF) and decode-and-forward (DF) : and seven neural methods spanning supervised learning (MLP, Hybrid SNR-adaptive relay), generative modeling (VAE, conditional GAN with WGAN-GP), and sequence architectures (Transformer, Mamba S6 selective state space model, Mamba-2 structured state space duality).

Evaluation covers six channel configurations: AWGN, Rayleigh fading, and Rician fading (K=3) in SISO topology, and 2×2 MIMO with Rayleigh fading under three equalization techniques : zero-forcing (ZF), minimum mean square error (MMSE), and successive interference cancellation (SIC). Primary experiments use BPSK modulation with Monte Carlo simulation (100,000 bits per SNR point, 95\% confidence intervals). Extension experiments cover QPSK, 16-QAM, and 16-PSK.

Key findings are as follows. Neural network relays show selective low-SNR gains (0--4 dB), most clearly on AWGN and Rician channels and in some MIMO configurations, but this advantage is not universal across all channels or architectures. Classical DF dominates at medium-to-high SNR (≥6 dB) with zero parameters. A complexity study reveals an inverted-U relationship between model size and performance: a minimal 169-parameter network matches models 100× larger, while an 11,201-parameter model overfits. A normalized 3,000-parameter comparison shows the performance gap between architectures narrows at equal scale, with VAE as the consistent underperformer : indicating that parameter count, not architecture, is the primary driver for this task. Mamba-2 SSD trains 35\% faster than Mamba S6 with comparable BER, exploiting parallel chunk computation. For MIMO, MMSE outperforms ZF and SIC provides further gains through interference cancellation.

The recommended deployment is a Hybrid relay combining neural processing at low SNR with DF at high SNR, achieving near-optimal performance with minimal overhead (169 parameters, \textasciitilde0.7 KB, under 3 seconds training). For higher-order modulations, a \textbf{joint 2D classifier} over all 16 constellation points eliminates the structural BER floor imposed by per-axis I/Q splitting: the top three 16-class variants (VAE, Transformer, MLP) reach near-zero BER at 20 dB, matching DF for the first time. A CSI injection study across 48 neural variants reveals modulation-dependence: CSI is detrimental for amplitude-carrying 16-QAM but beneficial for constant-envelope 16-PSK.

\textbf{Keywords:} Cooperative relay communication, deep learning, two-hop relay, MLP, Mamba S6, Mamba-2, Transformer, VAE, conditional GAN, MIMO equalization, 16-QAM, 16-PSK, CSI injection, bit error rate

\clearpage

%% ── Acknowledgments ──────────────────────────────────────────
\cleardoublepage
\phantomsection
\chapter*{Acknowledgments}
\addcontentsline{toc}{chapter}{Acknowledgments}
\markboth{Acknowledgments}{Acknowledgments}

I would like to dedicate this work to my wife \textit{Shelly Ines Zukerman Lavy} for supporting my graduate studies.I also would like to thank my supervisors \textit{Dr. Anatoly Khina} and \textit{Assoc. Prof.  Raja Giriys} at Tel Aviv University for their guidance and support throughout this research. 
This work was conducted in the School of Electrical and electronic Engineering, Tel Aviv University.


